\section{Related work}\label{sec:related-work}

\citet{GuoEtAl2026} formalized polynomial division with remainder,
Buchberger's criterion, reduced Gr\"obner bases, and results connecting
finite-variable and infinite-variable polynomial rings in Lean 4.  Therefore,
this article does not claim to be the first Lean formalization of Gr\"obner
basis theory.  The ideal-relative predicate \lean{IsGroebnerBasis} is a
standard definition also used in that development; this article does not
claim the definition itself as a new contribution.  In comparison, the
present development focuses on the explicit arbitrary-term reduction relation; the
equivalence between the Gr\"obner condition and Church--Rosser, local
confluence, confluence, reduction to zero, and uniqueness of normal forms;
and the well-founded nondeterministic Buchberger step relation.

\citet{ShenEtAl2026} developed tactics that send polynomial computations to
SageMath or SymPy and verify the returned certificates in Lean using a
computable polynomial representation.  Their system is designed for practical
automation on concrete inputs.  The formalization in this article does not call
an external computer algebra system: it proves the reduction theory and theorems about the abstract Buchberger
procedure inside Lean.  Conversely, it does not yet compute Gr\"obner bases on
concrete inputs.

In Isabelle/HOL, \citet{MaletzkyImmler2018} formalized Gr\"obner bases of
modules and executable implementations of Buchberger's and Faug\`ere's $F_4$
algorithms.  \citet{Maletzky2021} then formalized a generic executable family
of signature-based algorithms.  Those Isabelle developments include concrete
functions that can be evaluated, whereas \lean{BuchbergerStep} in this article
is a proposition-valued relation.  This difference motivates the future
separation between the relation \lean{BuchbergerStep} and an efficient
ordered representation for executable Lean code.
